\chapter{Variable Definitions and Data Pipeline}
\label{app:variables}
\label{app:diagnostics}

% TODO: Possible contents
%
% \section{Full Sample of Companies}
% - Table with ticker, CIK, company name, SIC industry, fiscal year-end
%
% \section{Variable Definitions}
% - Complete variable dictionary with XBRL tags and derivation formulae
% - Textual variable construction details
% - Divergence measure formal definition
%
% \section{Data Pipeline Documentation}
% - Description of SEC EDGAR extraction steps
% - Validation and cleaning procedures
% - Volatility and return computation methodology

\section{Filing Frequency Diagnostics}
\label{sec:freq_diagnostics}

The following tables supplement the discussion in Section~\ref{sec:annual_vs_quarterly}.

\begin{table}[htbp]
\centering
\caption{XBRL Financial Metric Coverage: Annual (10-K) vs.\ Quarterly (10-Q)}
\label{tab:xbrl_coverage}
\begin{tabular}{lrr}
\toprule
Metric & Annual (\%) & Quarterly (\%) \\
\midrule
assets\_current & 86.7 & 91.1 \\
eps\_diluted & 89.0 & 92.7 \\
liabilities\_current & 86.7 & 91.1 \\
net\_income & 93.2 & 97.3 \\
operating\_cash\_flow & 88.6 & 35.6 \\
operating\_income & 79.7 & 83.8 \\
stockholders\_equity & 98.7 & 94.9 \\
total\_assets & 89.7 & 93.3 \\
total\_liabilities & 66.5 & 69.7 \\
\midrule
Median periods per firm & 20 & 50 \\
Firms with data & 30 & 30 \\
\bottomrule
\end{tabular}
\begin{tablenotes}
\small
\item \textit{Note:} Coverage shows the percentage of firm-periods with non-missing
values for each metric. Annual data is extracted from 10-K filings; quarterly data
from 10-Q filings (Q1--Q3 only). Sample: 30 non-financial operating
companies from the top 200 by market capitalisation.
\end{tablenotes}
\end{table}


\begin{table}[htbp]
\centering
\caption{Annual vs.\ Quarterly Approach: Summary Comparison}
\label{tab:annual_vs_quarterly}
\begin{tabular}{p{5.5cm}p{4.5cm}p{4.5cm}}
\toprule
Criterion & Annual (10-K) & Quarterly (10-Q) \\
\midrule
Filing type & 10-K & 10-Q (Q1--Q3 only) \\
MD\&A availability & 93\% (Item 7) & 100\% (Part I, Item 2) \\
Risk Factors availability & 100\% (Item 1A) & $\sim$52\% substantive \\
Risk Factors content & Median $\sim$30{,}000 words & Median $\sim$9{,}000 words (substantive); 48\% boilerplate \\
Q4 financial data & Included in 10-K & Unreliable / absent \\
Explicit Q4 XBRL facts & N/A & 3--43\% of firm-years \\
Implied Q4 mismatch rate & N/A & $\sim$11\% of comparisons \\
Cash-flow quarterly coverage & N/A & 7\% of firm-years \\
NLP variable: $\Delta$Risk & Feasible & Limited (52\% substantive) \\
NLP variable: $\Delta$Sentiment & Feasible & Feasible (noisier) \\
NLP variable: TextSim & YoY (standard) & QoQ (unclear construct) \\
Literature alignment & Standard & Uncommon \\
Expected observations & $\sim$9{,}000 firm-years & $\sim$27{,}000 firm-quarters \\
\bottomrule
\end{tabular}
\begin{tablenotes}
\small
\item \textit{Note:} Summary based on diagnostic tests of 30 non-financial companies
(XBRL, Q4 tests) and 15 companies (textual tests). Risk Factors in 10-Q filings
are required by SEC Regulation S-K Item 1A only when there are ``material changes''
from the most recent 10-K; roughly half of companies include substantive risk
disclosures, while the rest provide only a boilerplate reference. See
Tables~\ref{tab:xbrl_coverage},~\ref{tab:q4_reliability},
and~\ref{tab:text_availability} for detailed results.
\end{tablenotes}
\end{table}


\begin{table}[htbp]
\centering
\caption{Per-Company 10-Q Textual Content (Most Recent 6 Filings)}
\label{tab:text_per_company}
\small
\begin{tabular}{lrrrrr}
\toprule
Ticker & Filings & MD\&A Found & MD\&A Med.\ Words & Risk Subst. & Risk Boilerplate \\
\midrule
AAPL & 6 & 6 & 2,781 & 3 & 3 \\
AMZN & 6 & 6 & 6,291 & 0 & 0 \\
AVGO & 6 & 6 & 4,778 & 6 & 0 \\
COST & 6 & 6 & 4,979 & 0 & 6 \\
GOOGL & 6 & 6 & 8,096 & 3 & 3 \\
JNJ & 6 & 6 & 7,587 & 0 & 6 \\
LLY & 6 & 6 & 4,994 & 0 & 6 \\
MA & 6 & 6 & 263 & 0 & 6 \\
META & 6 & 6 & 7,756 & 6 & 0 \\
MSFT & 6 & 0 & 189 & 0 & 6 \\
NVDA & 6 & 6 & 5,156 & 6 & 0 \\
TSLA & 6 & 6 & 5,582 & 2 & 4 \\
V & 6 & 6 & 5,043 & 0 & 6 \\
WMT & 6 & 6 & 7,142 & 0 & 6 \\
XOM & 6 & 6 & 6,197 & 6 & 0 \\
\bottomrule
\end{tabular}
\begin{tablenotes}
\small
\item \textit{Note:} ``MD\&A Found'' and ``Risk Subst.'' count filings with $>$200
substantive words. ``Risk Boilerplate'' counts filings where the risk-factors section
contains only a reference to the annual 10-K.
\end{tablenotes}
\end{table}


\section{XBRL Tag Provenance}
\label{sec:tag_provenance}

Table~\ref{tab:tag_provenance} reports, for each metric with multiple
candidate XBRL tags, the number and share of firms locked to each tag
by the per-firm tag-locking procedure described in
Section~\ref{sec:xbrl_tags}.

% TODO: Generate this table from data/diagnostics/tag_provenance.csv
% after re-running the pipeline.
\begin{table}[htbp]
  \centering
  \caption{Per-firm XBRL tag distribution}
  \label{tab:tag_provenance}
  \small
  \begin{tabular}{llrr}
    \toprule
    Metric & Locked XBRL tag & Firms & \% \\
    \midrule
    \multicolumn{4}{c}{\emph{To be populated after pipeline re-run}} \\
    \bottomrule
  \end{tabular}
  \begin{minipage}{0.9\linewidth}\vspace{4pt}\footnotesize
  \emph{Notes.} For each (firm, metric) pair, the tag with the highest
  within-firm year-count is selected; values from other tags are set to
  missing. See Section~\ref{sec:xbrl_tags} for the rationale.
  \end{minipage}
\end{table}
